%=======================================================================
%  CHAPTER 7 — APPLICATIONS OF AI  (14 hrs)
%=======================================================================
\section{Applications of AI}

{\color{sub}\footnotesize 14 hrs \gap$\cdot$\gap asked in \mk{all 41 papers}
\gap$\cdot$\gap 3 syllabus groups, 12 leaf topics \gap$\cdot$\gap 899 marks,
\mk{26.5\,\% of every mark ever set} \gap$\cdot$\gap the gate networks,
the Hopfield matrix, the forward propagation and the parse tree are worked in
\emph{Ch7 -- Applications of AI -- Solved Problems}.}

\lead{Read this first:}
\begin{itemize}
  \item \mk{This one chapter is a quarter of the paper.} It is bigger
        than chapters 1, 2, 3 and 5 put together, and \mk{no paper has
        ever skipped it}. Expect three or four questions from it every time.
  \item \mk{Four questions dominate, and each is nearly certain:}
  \begin{itemize}
    \item \textbf{The steps of NLP, with a block diagram} \tS{23} $-$ the
          most-asked question in the entire subject.
    \item \textbf{Perceptron, and design a net for AND / OR / XOR} \tS{17}.
    \item \textbf{Expert system architecture, with a block diagram} \tS{17}.
    \item \textbf{Machine vision: role, steps, applications} \tS{12}.
  \end{itemize}
  \item \mk{Three drawings are worth more than anything else here}:
        the five-box NLP pipeline, the four-box expert system, and a two-input
        perceptron. Practise all three until you can draw them in one minute each.
  \item \textbf{Then the neural networks}: Kohonen \tS{9}, Hopfield \tS{8},
        back-propagation \tS{8}, the neuron model and activation functions \tS{8},
        Hebbian learning \tF{4}.
  \item \mk{The examiner pairs across sub-topics}, so revise them
        together: \emph{machine vision $+$ back-propagation} (\yr{71 Ma},
        \yr{\textbf{\textit{72 Ash}}}), \emph{NLP challenges $+$ machine vision}
        (\yr{\textbf{\texttt{79 Bh}}}), \emph{declarative Vs procedural $+$ expert
        system architecture} (\yr{70 Ma}, \yr{72 Ma}, \yr{73 Ma}).
  \item \mk{Chapter 6's neural-network syllabus line is answered here.}
        The syllabus files neural networks under Machine Learning, but every paper
        sets them under Applications of AI, so perceptron, back-propagation, Hopfield,
        Kohonen and Hebbian learning are all in this chapter.
\end{itemize}

\hr

%-----------------------------------------------------------------------
\T{7.1 The artificial neuron: model, activation function, ANN Vs the brain}

\Q{What is an ANN? Explain the McCulloch-Pitts model with a diagram. What is the role of the activation function?}{\m{4}\m{3+4}\m{4+3}\m{1+7}\m{3+5}\m{8}\m{10}}
{\color{sub}\footnotesize \tS{8} \yr{\textbf{71 Bh} \m{3+5}, \textbf{81 Ch} \m{3+4},
82 Ka \m{4+3}, 75 Ba \m{8}, 76 Ba \m{10}, \textbf{\texttt{80 Bh}} \m{2+2+4},
\textbf{\texttt{79 Bh}} \m{1+7}, \textbf{69 Bh} \m{4}} \gap$\cdot$\gap
\yr{\textbf{71 Bh}} and \yr{\textbf{69 Bh}} want the \textbf{brain comparison},
\yr{\textbf{81 Ch}} and \yr{82 Ka} the \textbf{activation function},
\yr{75 Ba} and \yr{76 Ba} the \textbf{MP neuron on AND and XOR} (worked in the
companion), \yr{\textbf{\texttt{80 Bh}}} the \textbf{mathematical model and what
training means}}

\sbsr{0.55}{%
\lead{Definition and the mathematical model}
\begin{itemize}
  \item \textbf{An artificial neural network is a parallel computational system of many
        simple processing elements (neurons), connected in a specific way, that learns a
        task from data rather than from a program.}
  \item Modelled on the brain: massively parallel, and \mk{tolerant of
        noise and of the failure of individual units}.
  \item \textbf{One neuron does three things}, and that is the whole model:
  \begin{itemize}
    \item \textbf{1. Weight} each input. $w_1 x_1$ up to $w_n x_n$. The weight is the
          strength of that connection, and \mk{the weights are what
          the network learns}.
    \item \textbf{2. Sum} them, plus a \textbf{bias} $b$ (equivalently, subtract a
          \textbf{threshold} $\theta$):
          \[ net \;=\; \sum_{i=1}^{n} w_i x_i \;+\; b \]
    \item \textbf{3. Apply an activation function} $f$ to that sum:
          $y = f(net)$.
  \end{itemize}
\end{itemize}}{%
\figT{c7_mcp.png}
{\color{sub}\footnotesize The McCulloch-Pitts neuron and its firing condition. Weight,
sum, threshold: the same three operations every neuron in this chapter performs.}}

\sbsr{0.55}{%
\begin{itemize}
  \item \textbf{The five elements to name} if the question says ``elements of a neural
        network'': \textbf{inputs}, \textbf{weights}, \textbf{threshold or bias},
        \textbf{activation function}, \textbf{learning rate} $\alpha$.
  \item \mk{What ``training a network'' means} (\yr{\textbf{\texttt{80 Bh}}}
        \m{2}): presenting examples and adjusting the weights and biases so the output
        moves towards the target. \textbf{Nothing else in the network changes.}
  \item \textbf{Learning rate} $\alpha$ sets the step size. Too large and the weights
        oscillate and never settle; too small and training takes forever. Typical
        $0.01$ to $0.5$.
\end{itemize}}{%
\lead{The McCulloch-Pitts neuron (1943)}
{\color{sub}\footnotesize the first neuron model, and the one four papers name.}
\begin{itemize}
  \item Output is \textbf{binary}: 1 (fires) or 0. Nothing in between.
  \item Every input is weighted; the neuron \textbf{fires when the weighted sum exceeds
        the threshold} $T$:
        \[ x_1w_1 + x_2w_2 + x_3w_3 > T \;\Rightarrow\; y=1 \]
  \item \textbf{Excitatory} inputs carry positive weights, \textbf{inhibitory} inputs
        negative ones, and in the original model \mk{one active
        inhibitory input vetoes the neuron outright}.
  \item \mk{The weights are fixed by the designer, not learnt.} That
        is the difference from a perceptron, and it is the sentence the examiner is
        looking for.
  \item It can realise \textbf{AND, OR and NOT} $-$ so any Boolean function, if you are
        allowed several layers $-$ but \mk{one MP neuron cannot do
        XOR}. Both worked in the companion.
\end{itemize}}

\sbsr{0.42}{%
\figT{c7_activation.png}}{%
\lead{The activation function, and how to choose one \gap{\color{sub}\footnotesize(\yr{\textbf{81 Ch}} \m{3}, \yr{82 Ka} \m{3})}}
\begin{itemize}
  \item \textbf{What it is for}, three answers, all worth marks:
  \begin{itemize}
    \item \textbf{Decides whether the neuron fires}, by mapping an unbounded sum onto a
          bounded output.
    \item \mk{Introduces non-linearity}. Without it, a stack of
          layers collapses into one linear map, and the network can only ever learn a
          straight line $-$ this is the answer to \yr{\textbf{\texttt{79 Bh}}}'s
          \emph{``how can non-linearity be modelled through an ANN''}.
    \item Must be \textbf{differentiable} for gradient descent, which is why
          back-propagation needs sigmoid or tanh, not a step.
  \end{itemize}
  \item \textbf{Choosing one for a project} $-$ answer with the rule, then an example:
  \begin{itemize}
    \item \textbf{Binary classification output} $\rightarrow$ \textbf{sigmoid}, because
          it gives a value in $(0,1)$ readable as a probability.
    \item \textbf{Multi-class output} $\rightarrow$ \textbf{softmax}.
    \item \textbf{Regression output} $\rightarrow$ \textbf{linear}, so the value is
          unbounded.
    \item \textbf{Hidden layers of a deep net} $\rightarrow$ \textbf{ReLU}: cheap, and
          it does not saturate, so it does not kill the gradient. Its cost is
          \mk{dead units} $-$ a neuron whose input stays negative
          has zero gradient and never recovers, which is exactly what happens in the
          numerical worked in the companion.
    \item \textbf{Zero-centred hidden layer} $\rightarrow$ \textbf{tanh}, which
          converges faster than sigmoid for that reason.
    \item \textbf{Hard threshold} (step, signum) $\rightarrow$ only for a perceptron or
          an MP neuron, never where gradients are needed.
  \end{itemize}
\end{itemize}}

\sbs{%
\lead{ANN Vs the human brain \gap{\color{sub}\footnotesize(\yr{\textbf{71 Bh}} \m{5}, \yr{\textbf{69 Bh}} \m{4})}}
\begin{tabularx}{\linewidth}{@{}L{2.2cm}XX@{}}
\hrow \thead{} & \thead{Human brain} & \thead{ANN} \\
Units & about \textbf{86 billion} neurons & a few tens to a few billion \\
Connections & about \textbf{100 trillion} synapses & set by the architecture \\
Speed & slow: a neuron fires at a few hundred Hz & fast: gigahertz \\
Processing & \textbf{massively parallel}, all at once & parallel in principle, usually
  serial on one machine \\
Memory & about 2.5 petabytes, \textbf{content addressable} & limited, address based \\
Fault tolerance & \textbf{very high}: neurons die daily and nothing is lost & high
  within a net, none if the machine fails \\
Learning & continuous, from few examples & needs many examples and retraining \\
Power & about \textbf{20 W} & kilowatts to train \\
\end{tabularx}
\begin{itemize}
  \item \mk{The contrast worth quoting}: for vision the brain is worth a
        thousand supercomputers; for arithmetic a pocket calculator is worth ten brains.
        Structure, not speed, is what differs.
\end{itemize}}{%
\lead{Network structures}
\figT{c7_mlp.png}
\begin{itemize}
  \item \textbf{Feed-forward} $-$ signals travel one way, input to output, with
        \mk{no loops}. A layer's output never affects that same layer.
  \begin{itemize}
    \item \textbf{Single layer}: input straight to output. Only linearly separable
          problems (\S7.3).
    \item \textbf{Multilayer (MLP)}: one or more \textbf{hidden layers} between them.
          \mk{One hidden layer is enough to represent any Boolean
          function}, and with enough units, any continuous function.
  \end{itemize}
  \item \textbf{Feedback (recurrent)} $-$ loops allowed, so the network has
        \mk{state, and therefore memory}. More powerful, harder to
        train. Hopfield (\S7.5) is the standard example.
  \item \textbf{Competitive} $-$ output units compete and only the winner fires.
        Kohonen (\S7.6).
\end{itemize}}

%-----------------------------------------------------------------------
\T{7.2 Perceptron, linear separability and the XOR problem}

\Q{What is a perceptron? Design a neural network that acts as an AND / OR / XOR gate.}{\m{4}\m{6}\m{7}\m{10}\m{1+7}\m{2+6}\m{1+8}}
{\color{sub}\footnotesize \tS{17} \yr{72 Ma \m{2+7}, 73 Ma \m{1+7},
\textit{74 Ma} \m{10}, \textbf{\textit{74 Bh}} \m{1+8}, \textbf{\texttt{82 Bh}} \m{2+6},
\textbf{79 Ch} \m{2+6}, 79 Ash \m{1+8}, \textbf{77 Ch} \m{2+6}, 79 Jth \m{1+6},
\textbf{\textit{73 Bh}} \m{2+6}, 78 Po \m{2+2+4}, \textbf{\textit{77 Ch}} \m{2+6},
\textit{72 Ma} \m{7}, \textbf{\textit{76 Bh}} \m{6}, \textbf{\texttt{81 Bh}} \m{2+6},
69 Po \m{4}, \textbf{72 Ash} \m{2+4+4}} \gap$\cdot$\gap
\mk{second most-asked question in the subject}, and eight of the
seventeen name \textbf{XOR} specifically \gap$\cdot$\gap all four gates are worked in
the companion}

\sbsr{0.55}{%
\lead{The perceptron}
\begin{itemize}
  \item \textbf{A perceptron is a single artificial neuron with adjustable weights, a
        bias and a threshold activation, trained by an iterative rule that is guaranteed
        to converge when the classes are linearly separable.} Rosenblatt, \textbf{1958}.
  \item \mk{The difference from a McCulloch-Pitts neuron}: the MP
        neuron's weights are chosen by hand, the perceptron's are \emph{learnt}.
  \item \textbf{Activation}, with threshold $\theta$:
        \[ f(net) = \begin{cases} 1 & net > \theta \\ 0 & -\theta \le net \le \theta \\
           -1 & net < -\theta \end{cases} \]
  \item \textbf{The perceptron learning rule.} For each training pair, compute
        $y$; if $y \neq t$, update, otherwise change nothing:
        \[ w_i(new) = w_i(old) + \alpha\, t\, x_i \qquad
           b(new) = b(old) + \alpha\, t \]
  \item \mk{Weights change only when there is an error.} That single
        sentence separates the perceptron rule from Hebbian learning (\S7.4), which
        updates on \emph{every} pattern.
  \item \textbf{Convergence theorem}: if a separating line exists, the rule finds one in
        a finite number of steps, whatever the starting weights.
\end{itemize}}{%
\figT{c7_perceptron.png}
\lead{Limitations \gap{\color{sub}\footnotesize(\yr{\textbf{\textit{76 Bh}}} \m{6} asks precisely this)}}
\begin{itemize}
  \item Learns \mk{only linearly separable} classes.
  \item If the data is not separable the algorithm does not converge and its behaviour
        is not defined $-$ it cycles.
  \item If it \emph{is} separable there are usually \textbf{many} valid answers, and the
        rule gives no way to tell whether the one it found is a good one (no margin, no
        confidence). That is what SVMs later fixed.
\end{itemize}}

\sbs{%
\lead{Linear separability, and why it decides everything}
\begin{itemize}
  \item A single neuron computes $w_1x_1 + w_2x_2 + b$ and thresholds it. The set of
        inputs where that sum is zero is \mk{a straight line} (a
        hyperplane in $n$ dimensions). Everything on one side outputs 1, everything on
        the other outputs 0.
  \item So \mk{a perceptron can learn a function if and only if one
        straight line separates its 1s from its 0s}. That is the whole of
        \yr{\textbf{\texttt{81 Bh}}}'s \emph{``perceptron works as a linear classifier''}.
  \item Plot the four corners of the unit square:
  \begin{itemize}
    \item \textbf{AND} $-$ only $(1,1)$ is 1. One line cuts it off. \checkmark
    \item \textbf{OR} $-$ only $(0,0)$ is 0. One line cuts it off. \checkmark
    \item \textbf{XOR} $-$ the two 1s sit at \emph{opposite} corners, with the two 0s at
          the other opposite corners. \mk{No straight line separates
          them.} \xmark
  \end{itemize}
  \item \mk{Draw that square.} Marks are given for the picture, and
        it takes ten seconds.
\end{itemize}}{%
\lead{The algebraic proof, which is what a ``justify'' question wants}
\begin{itemize}
  \item Suppose weights $w_1, w_2$ and bias $b$ realise XOR with the rule
        \emph{output 1 iff $w_1x_1 + w_2x_2 + b > 0$}. The four rows give:
  \begin{itemize}
    \item $(0,0) \rightarrow 0$: \quad $b \le 0$ \hfill (i)
    \item $(0,1) \rightarrow 1$: \quad $w_2 + b > 0$ \hfill (ii)
    \item $(1,0) \rightarrow 1$: \quad $w_1 + b > 0$ \hfill (iii)
    \item $(1,1) \rightarrow 0$: \quad $w_1 + w_2 + b \le 0$ \hfill (iv)
  \end{itemize}
  \item Add (ii) and (iii): $w_1 + w_2 + 2b > 0$, so $w_1 + w_2 > -2b$.
  \item From (iv): $w_1 + w_2 \le -b$. Together: $-2b < -b$, i.e.
        \mk{$b > 0$}, which contradicts (i).
  \item \textbf{Therefore no single perceptron can implement XOR.} $\blacksquare$
  \item \mk{The fix}: a \textbf{multilayer} perceptron. One hidden
        layer computes OR and NAND, and the output neuron ANDs them, so the two hidden
        units bend the decision boundary into two lines. Built and checked row by row in
        the companion.
\end{itemize}}

%-----------------------------------------------------------------------
\T{7.3 Learning in a neural network: Hebbian, delta and back-propagation}

\Q{What is the back-propagation algorithm? When do we need it? Explain with an example.}{\m{4}\m{2+2+4}\m{3+5}\m{4+4}\m{2+6}\m{2+4+4}}
{\color{sub}\footnotesize \tS{8} \yr{\textbf{70 Bh} \m{4+4},
\textbf{72 Ash} \m{2+4+4}, 70 Ma \m{3+5}, \textbf{\texttt{80 Bh}} \m{2+2+4},
\textbf{74 Bh} \m{4}, \textbf{79 Ch} \m{2+6}, 71 Ma \m{3+5},
\textbf{\textit{72 Ash}} \m{3+5}} \gap$\cdot$\gap \yr{71 Ma} and
\yr{\textbf{\textit{72 Ash}}} attach it to \textbf{machine vision} (\S7.9), so revise the
two together \gap$\cdot$\gap one full step is computed in the companion}

\sbsr{0.53}{%
\lead{The three learning rules, in one place}
\begin{tabularx}{\linewidth}{@{}L{2.0cm}XX@{}}
\hrow \thead{Rule} & \thead{Update} & \thead{Updates when} \\
\textbf{Hebb} & $w_i \mathrel{+}= x_i\,t$ & \mk{every} pattern \\
\textbf{Perceptron} & $w_i \mathrel{+}= \alpha\,t\,x_i$ & only on an \textbf{error} \\
\textbf{Delta (LMS, Widrow-Hoff)} & $w_i \mathrel{+}= \alpha\,(t - y_{in})\,x_i$ &
  every pattern, by \textbf{how wrong} it was \\
\end{tabularx}
\begin{itemize}
  \item The delta rule uses the \emph{net input} $y_{in}$, not the thresholded output, so
        it measures the size of the error and not just its sign. That is what makes it
        differentiable, and therefore what back-propagation generalises.
  \item \textbf{Adaline} (Adaptive Linear Neuron, Widrow and Hoff 1960) is a single unit
        trained by the delta rule: identity activation while training, threshold
        afterwards. \mk{On the syllabus, never yet examined}, so know
        it in four lines and no more.
\end{itemize}}{%
\figT{c7_bp_xor.png}
\lead{Back-propagation}
\begin{itemize}
  \item \textbf{Short for ``backward propagation of errors''}: a supervised algorithm
        that trains a \textbf{multilayer} network by gradient descent, pushing the output
        error backwards through the layers to get a gradient for every weight.
  \item \mk{When do we need it} (\yr{\textbf{79 Ch}} \m{2}): as soon
        as there is a \textbf{hidden layer}. A hidden unit has no target of its own, so
        the perceptron rule cannot say what its error is. Back-propagation manufactures
        one from the chain rule. No hidden layer, no XOR (\S7.2), so no useful network.
  \item \textbf{Why sigmoid}: $f(x)=\frac{1}{1+e^{-x}}$ has the derivative
        $f'(x) = f(x)\,(1-f(x))$, so the forward pass already computed everything the
        backward pass needs. That is the only reason it is the default in every
        textbook example.
\end{itemize}}

\sbsr{0.53}{%
\lead{Hebbian learning \gap{\color{sub}\footnotesize\tF{4} \yr{\textbf{73 Bh} \m{3+7}, \texttt{81 Ba} \m{2+6}, 78 Ba \m{1+7}, \textbf{78 Ch} \m{6+3}}}}
\begin{itemize}
  \item Hebb's postulate, 1949: \emph{when cell A repeatedly takes part in firing cell B,
        the connection from A to B is strengthened}. Usually quoted as
        \mk{``neurons that fire together wire together''}.
  \item \textbf{Algorithm} (four steps, and this is the whole answer):
  \begin{itemize}
    \item \textbf{Step 0.} Set all weights and the bias to \textbf{0}.
    \item \textbf{Step 1.} For a training pair $(s, t)$, set the inputs $x_i = s_i$.
    \item \textbf{Step 2.} Set the output to the target, $y = t$.
    \item \textbf{Step 3.} $w_i(new) = w_i(old) + x_i y$ and $b(new) = b(old) + y$.
    \item Repeat for every pair. \mk{One pass. No iteration, no
          learning rate, no error term.}
  \end{itemize}
  \item \mk{Is Hebbian learning supervised?} (\yr{\textbf{78 Ch}}
        \m{3}) $-$ \textbf{yes, as used here}: the update needs the target $t$, so the
        network is told the right answer. Hebb's biological postulate is unsupervised;
        the \emph{Hebb net} for pattern association is not. Say both halves.
  \item It needs \textbf{bipolar} data. With binary inputs a 0 contributes nothing to any
        weight, so the AND net cannot be learnt at all $-$ a real trap, and the reason
        every worked example uses $\pm 1$.
\end{itemize}}{%
\lead{The back-propagation algorithm}
\begin{itemize}
  \item \textbf{Six steps:}
  \begin{itemize}
    \item \textbf{1.} Initialise all weights and biases to small random values, pick
          $\eta$.
    \item \textbf{2. Forward pass}: layer by layer, compute
          $net_j = \sum w_{ij} x_i + b_j$ and $h_j = f(net_j)$ to the output $y$.
    \item \textbf{3. Error} at the output, $E = \tfrac12 (y - t)^2$.
    \item \textbf{4. Backward pass}: output delta
          $\delta_{out} = (y - t)\,y\,(1-y)$; hidden delta
          $\delta_j = h_j(1-h_j)\sum_k \delta_k w_{jk}$.
    \item \textbf{5. Update} every weight, $w \leftarrow w - \eta\,\delta\,(\text{input
          to that weight})$, and every bias, $b \leftarrow b - \eta\,\delta$.
    \item \textbf{6.} Repeat over the training set until $E$ is small enough or the
          epoch limit is reached.
  \end{itemize}
  \item \textbf{Problems to name}: slow; \mk{gets stuck in local
        minima}; \textbf{vanishing gradient} through many sigmoid layers; sensitive
        to $\eta$ and to the initial weights; can overfit.
\end{itemize}}

%-----------------------------------------------------------------------
\T{7.4 Hopfield network}

\Q{What is a Hopfield network? Explain its features and all its steps with an example.}{\m{3}\m{4}\m{6}\m{8}\m{3+5}}
{\color{sub}\footnotesize \tS{8} \yr{\textbf{69 Bh} \m{8}, \textbf{75 Bh} \m{8},
\textbf{\textit{71 Bh}} \m{8}, 80 Ash \m{3+5}, \textbf{\texttt{82 Bh}} \m{6},
\textbf{80 Ch} \m{4}, \texttt{82 Ba} \m{3}, 71 Ma \m{3}} \gap$\cdot$\gap
\yr{\textbf{\texttt{82 Bh}}} wants the \textbf{weight matrix and a recall computed}, which
is the companion \gap$\cdot$\gap \yr{80 Ash} and \yr{71 Ma} want it
\textbf{compared} $-$ with a feed-forward net and with Kohonen}

\sbsr{0.55}{%
\begin{itemize}
  \item \textbf{A Hopfield network is a single-layer, fully connected, recurrent network
        of binary units with symmetric weights, which settles from any starting state
        into the nearest stored pattern.} Hopfield, \textbf{1982}.
  \item \textbf{Structure}, and every word of it earns a mark:
  \begin{itemize}
    \item One layer. Every unit is connected to \textbf{every other} unit, and each
          unit's output feeds back as an input.
    \item Weights are \textbf{symmetric}, $w_{ij} = w_{ji}$, with
          \mk{no self connection}, $w_{ii}=0$.
    \item Units are \textbf{binary} $\{0,1\}$ or \textbf{bipolar} $\{-1,+1\}$.
    \item \mk{There is no separate output layer}: the same units
          take the input and hold the answer.
  \end{itemize}
  \item \textbf{Hopfield Vs feed-forward} (\yr{80 Ash} \m{5}): recurrent Vs one-way;
        symmetric weights Vs arbitrary; one layer Vs many; associative memory Vs function
        approximation; converges to a stored pattern Vs computes an output in one pass;
        trained in \textbf{one shot} Vs trained by many epochs of back-propagation.
\end{itemize}}{%
\figT{c7_hopfield_states.png}
{\color{sub}\footnotesize Four stable states of one small net. Supply any part of one and
the network falls into it.}
\figT{c7_hopfield_energy.png}
{\color{sub}\footnotesize Every update walks downhill on this surface, so the network
always stops somewhere. The valleys are the stored patterns.}}

\sbsr{0.55}{%
\begin{itemize}
  \item \textbf{The four features to list} (\yr{\textbf{\textit{71 Bh}}} asks for them by
        name):
  \begin{itemize}
    \item \textbf{Content-addressable memory} $-$ retrieve a whole pattern by supplying
          \emph{part} of it. This is the headline property.
    \item \textbf{Distributed representation} $-$ a pattern is stored in all the weights
          at once, not in one cell.
    \item \textbf{Asynchronous control} $-$ units update one at a time, in any order, with
          no central clock.
    \item \textbf{Fault tolerance} $-$ damage a few weights and recall degrades
          gracefully rather than failing.
  \end{itemize}
  \item \textbf{Energy.} Each state has an energy
        \[ E = -\tfrac12 \sum_{i \neq j} w_{ij}s_i s_j - \sum_i \theta_i s_i \]
        and every update either lowers it or leaves it alone. So the network
        \mk{always converges}, to a local minimum of $E$. The stored
        patterns are those minima.
  \item \textbf{Capacity}: about $p \le 0.15N$ patterns for $N$ units. Store more and the
        patterns interfere; the net then settles into \textbf{spurious states} that were
        never taught.
\end{itemize}}{%
\lead{The algorithm}
\begin{itemize}
  \item \textbf{Training} (one pass, no iteration). For $P$ bipolar patterns:
        \[ w_{ij} \;=\; \sum_{p=1}^{P} s_i^{(p)} s_j^{(p)} \quad (i \neq j),
           \qquad w_{ii}=0 \]
        i.e. sum the outer products and zero the diagonal.
  \item \textbf{Recall / testing}, given a noisy input $x$:
  \begin{itemize}
    \item \textbf{1.} Set the state $y = x$.
    \item \textbf{2.} Pick one unit $i$ at random and compute
          $y_{in,i} = x_i + \sum_j y_j w_{ji}$.
    \item \textbf{3.} Apply the threshold: $y_i = 1$ if $y_{in,i} > \theta_i$,
          unchanged if equal, $0$ (or $-1$) if less.
    \item \textbf{4.} \mk{Broadcast the new $y_i$ immediately} $-$
          the next unit uses it. This is what \emph{asynchronous} means, and updating
          all units at once is the commonest mistake.
    \item \textbf{5.} Repeat until no unit changes. The state is then stable.
  \end{itemize}
  \item \mk{Worked in full in the companion}: three patterns stored,
        the weight matrix built, and a corrupted pattern recalled unit by unit.
\end{itemize}}

%-----------------------------------------------------------------------
\T{7.5 Kohonen network and the self-organizing map}

\Q{What is a self-organizing map? Explain all the steps involved with an example.}{\m{3}\m{4}\m{5}\m{2+5}}
{\color{sub}\footnotesize \tS{9} \yr{\textbf{74 Bh} \m{2+5},
\textbf{\texttt{80 Bh}} \m{4}, \textbf{\texttt{79 Bh}} \m{4}, \textbf{75 Bh} \m{4},
\textit{74 Ma} \m{4}, \textbf{\textit{72 Ash}} \m{4}, 79 Jth \m{5},
\textbf{\texttt{81 Bh}} \m{3}, 71 Ma \m{3}} \gap$\cdot$\gap
\mk{seven of the nine ask it as a short note}, so the answer is a
definition, the algorithm as five steps, and one application \gap$\cdot$\gap
\yr{71 Ma} wants it \textbf{compared with Hopfield}}

\sbsr{0.55}{%
\begin{itemize}
  \item \textbf{A self-organizing map (Kohonen network) is an unsupervised, competitive
        neural network that maps high-dimensional input onto a low-dimensional grid of
        neurons while preserving the topology of the data.} Kohonen, \textbf{1982}.
  \item \textbf{Structure}: an input layer, and \mk{a grid of output
        neurons} (usually 2-D). Fully connected between them and
        \mk{nothing else} $-$ no hidden layer, no bias, no activation
        function.
  \item \textbf{Competitive learning}: for each input, exactly one output neuron
        \mk{wins}, the one whose weight vector is nearest the input.
        It and its \emph{neighbours} are pulled towards that input; everyone else is
        left alone.
  \item \mk{Topology preserving} is the property to name: inputs that
        are similar end up firing neurons that are close together on the grid. That is
        what makes an SOM a visualisation tool and not just a clustering algorithm.
  \item \textbf{Used for}: clustering, dimensionality reduction, data visualisation
        (the poverty map is the classic), feature extraction, image compression.
\end{itemize}}{%
\figT{c7_som.png}
{\color{sub}\footnotesize The grid of neurons, and the same net drawn as connections from
the input vector to every output neuron with one marked \emph{winner}.}}

\sbsr{0.50}{%
\lead{Hopfield Vs Kohonen \gap{\color{sub}\footnotesize(\yr{71 Ma} \m{3})}}
\begin{tabularx}{\linewidth}{@{}L{1.7cm}XX@{}}
\hrow \thead{} & \thead{Hopfield} & \thead{Kohonen} \\
Learning & supervised (patterns given) & \textbf{unsupervised} \\
Structure & one layer, fully recurrent & input layer $\rightarrow$ output grid,
  feed-forward \\
Weights & symmetric, $w_{ii}=0$ & one vector per output neuron \\
Rule & Hebbian outer product, one shot & competitive, winner takes most \\
Output & a stored pattern & the winning neuron's position \\
Used for & associative memory, optimization & clustering, visualisation \\
\end{tabularx}
\begin{itemize}
  \item \checkmark Easy to interpret; handles large, complex data. \xmark
        Computationally expensive; the grid size and the initial weights are guesses;
        it assumes nearby points should behave alike.
\end{itemize}}{%
\lead{The algorithm, five steps}
\begin{itemize}
  \item \textbf{Step 1. Initialise} every neuron's weight vector $w_j$ randomly, and set
        the learning rate $\alpha$ and the neighbourhood radius $\sigma$.
  \item \textbf{Step 2. Present} an input vector $x$.
  \item \textbf{Step 3. Find the winner} (the Best Matching Unit), the neuron whose
        weight vector is nearest:
        \[ D(j) = \sum_{i} (x_i - w_{ij})^2, \qquad
           j^{*} = \arg\min_j D(j) \]
        {\color{sub}\footnotesize the \emph{squared} Euclidean distance is enough $-$
        the square root does not change which one is smallest, and skipping it saves
        time in an exam.}
  \item \textbf{Step 4. Update the winner and its neighbours}, moving them towards the
        input:
        \[ w_{ij}(new) = w_{ij}(old) + \alpha\,\bigl(x_i - w_{ij}(old)\bigr) \]
        for every $j$ inside the neighbourhood of $j^{*}$.
  \item \textbf{Step 5. Shrink} $\alpha$ and $\sigma$, and repeat from step 2 until they
        reach their floor or the map stops moving.
  \item \textbf{Two phases worth naming}: an early \textbf{ordering} phase with a large
        radius that lays the map out, and a long \textbf{convergence} phase with a small
        radius that fine-tunes it.
\end{itemize}}

%-----------------------------------------------------------------------
\T{7.6 Expert systems: what they are, and the architecture}

\Q{What is an expert system? Explain its general architecture with a block diagram.}{\m{4}\m{6}\m{8}\m{2+4}\m{1+7}\m{2+5}\m{3+5}}
{\color{sub}\footnotesize \tS{17} \yr{69 Po \m{8}, 71 Ma \m{2+4}, 76 Ba \m{1+7},
\textbf{\textit{71 Bh}} \m{8}, \texttt{80 Ba} \m{2+5}, \texttt{81 Ba} \m{3+5},
\textbf{77 Ch} \m{6+2}, \textbf{78 Ch} \m{1+4+3}, \textbf{\textit{76 Bh}} \m{5+2+2},
70 Ma \m{2+6}, 72 Ma \m{3+6}, 73 Ma \m{3+5}, 75 Ba \m{2+6}, \textbf{80 Ch} \m{6},
\texttt{82 Ba} \m{8}, \textbf{71 Bh} \m{2+6}, \textbf{74 Bh} \m{2+5}}
\gap$\cdot$\gap \mk{joint most-asked question in the subject}
\gap$\cdot$\gap three papers dress it as a scenario $-$ \emph{disease diagnosis}
(\yr{\textbf{80 Ch}}), \emph{a disaster-prone area of Nepal} (\yr{\texttt{82 Ba}}),
\emph{create a scenario where we require one} (\yr{\textbf{79 Ch}}) $-$ and want the
\emph{same} block diagram with the domain renamed}

\sbsr{0.52}{%
\lead{Definition}
\begin{itemize}
  \item \textbf{An expert system is a computer program that simulates the
        decision-making of a human expert in one narrow domain, using a body of
        high-quality specific knowledge rather than a general algorithm.}
  \item \mk{Narrow is the operative word.} An expert system is expert
        at one thing and useless at everything else.
  \item Knowledge comes from \textbf{human specialists}, and also from texts, journals
        and databases.
\end{itemize}}{%
\figT{c7_es_block.png}
{\color{sub}\footnotesize \mk{Draw this.} Four boxes and one
double-headed arrow to the user is a full answer to the block-diagram half of every one
of the seventeen papers.}}

\sbsr{0.52}{%
\begin{itemize}
  \item \textbf{Expert system shell} (\yr{82 Ka} \m{2}): an expert system with the
        \mk{knowledge base emptied out} $-$ inference engine, user
        interface and explanation facility, ready for a new domain's rules. It is what
        turns building an expert system into writing rules instead of writing code.
        Eg EMYCIN, the shell left when MYCIN's medical rules were removed.
  \item \textbf{When is a problem suited to one} (\yr{82 Ka} \m{3}): the domain is
        \textbf{narrow and well bounded}; genuine \textbf{human experts exist} and can
        explain their reasoning; the task is \textbf{symbolic and heuristic}, not
        numerical; it takes an expert minutes to hours, not seconds or months; and no
        ordinary algorithm solves it.
\end{itemize}}{%
\lead{The components}
\begin{itemize}
  \item \textbf{Knowledge base} $-$ the domain knowledge, almost always as
        \textbf{IF--THEN production rules}: \emph{IF the battery is dead THEN the car
        will not start}. This is what makes it an expert system rather than a program.
  \item \textbf{Working memory} $-$ the facts about \emph{this} case, built up during the
        consultation. Sometimes called the database or the fact base.
  \item \textbf{Inference engine} $-$ matches rules against working memory and fires
        them. Also called the control structure or rule interpreter.
  \item \textbf{User interface} $-$ two-way: it takes the user's queries and answers, and
        it presents conclusions.
  \item \textbf{Explanation facility} $-$ answers \emph{why are you asking} and
        \emph{how did you conclude that}. \mk{Name it: it is the one
        component most answers forget}, and it is what makes an expert system
        acceptable to a professional.
  \item \textbf{Knowledge acquisition subsystem} $-$ how new rules get in (\S7.7).
\end{itemize}}

\sbsr{0.52}{%
\lead{How the inference engine works \gap{\color{sub}\footnotesize(\yr{\texttt{81 Ba}} \m{5}, \yr{\textbf{\texttt{82 Bh}}} \m{5})}}
\begin{itemize}
  \item It applies one rule of logic, \textbf{modus ponens}: if \emph{A $\rightarrow$ B}
        is in the knowledge base and \emph{A} is in working memory, add \emph{B}.
  \item \textbf{Two directions}, both from Chapter 4:
  \begin{itemize}
    \item \textbf{Forward chaining (data driven)} $-$ start from the facts, fire whatever
          matches, add the conclusions, repeat. Use it when there are few facts and many
          possible conclusions. \emph{Monitoring, planning, design.}
    \item \textbf{Backward chaining (goal driven)} $-$ start from a hypothesis and look
          for rules that would prove it, asking the user for what is missing. Use it when
          there are few plausible conclusions. \emph{Diagnosis.} MYCIN works this way,
          and it is why the system asks you questions.
  \end{itemize}
  \item \textbf{The recognize--select--act cycle}, which is the answer if the question
        says \emph{how is inference performed}:
  \begin{itemize}
    \item \textbf{1. Match} $-$ compare every rule's condition against working memory,
          producing a \textbf{conflict set} of rules that could fire.
    \item \textbf{2. Conflict resolution} $-$ choose \emph{one} of them, by specificity,
          recency, rule priority or a \textbf{meta-rule}.
    \item \textbf{3. Act} $-$ fire it, updating working memory. Repeat until the goal is
          reached or nothing matches.
  \end{itemize}
  \item \textbf{Rule types to name}: \textbf{relationship} (\emph{IF the battery is dead
        THEN the car will not start}), \textbf{recommendation} (\emph{... THEN take a
        cab}), \textbf{directive} (\emph{... AND the fuel system is ok THEN check the
        electrical system}), \textbf{heuristic}, and \textbf{meta-rules}, which are rules
        about which rules to use.
\end{itemize}}{%
\figT{c7_es_arch.png}
{\color{sub}\footnotesize The fuller picture, worth drawing when the question is worth 6
or more: the dashed line splits the \textbf{consultation environment} the end user sees
from the \textbf{development environment} the knowledge engineer works in.}
\lead{Two real systems to name}
\begin{itemize}
  \item \textbf{MYCIN} (Stanford, 1970s) $-$ diagnoses blood infections and recommends
        antibiotics from about \textbf{500 production rules}. Uses
        \textbf{backward chaining}, and will explain its reasoning on request. Matched
        human specialists.
  \item \textbf{DENDRAL} (Stanford, late 1960s) $-$ the \textbf{first} expert system.
        Infers molecular structure from mass spectra using chemists' heuristics.
        \textbf{Forward chaining}.
  \item Others worth one word each: \textbf{XCON} (configuring VAX computers),
        \textbf{PROSPECTOR} (mineral exploration), \textbf{CADUCEUS} (internal medicine).
\end{itemize}}

%-----------------------------------------------------------------------
\T{7.7 Expert system development, and knowledge acquisition}

\Q{Explain the stages of expert system development, with an example.}{\m{6}\m{8}\m{2+5}\m{2+6}\m{4+4}}
{\color{sub}\footnotesize \tS{8} \yr{\textbf{69 Bh} \m{8}, \textbf{70 Bh} \m{4+4},
\textbf{73 Bh} \m{2+6}, 75 Ba \m{2+6}, \textbf{79 Ch} \m{2+6}, \textbf{80 Ch} \m{6},
\texttt{82 Ba} \m{8}, \texttt{80 Ba} \m{2+5}} \gap$\cdot$\gap
\mk{always paired with the block diagram}, so answer both halves
\gap$\cdot$\gap \yr{\textbf{80 Ch}} casts you as the \textbf{knowledge engineer} for a
disease-diagnosis system, \yr{\texttt{82 Ba}} in a \textbf{disaster-prone area of Nepal}}

\sbsr{0.52}{%
\lead{The six stages}
\begin{itemize}
  \item \textbf{1. Identify the problem domain.} Is the problem suitable at all? Are
        there real experts, and will they cooperate? Is it cost-effective?
  \item \textbf{2. Design the system.} Choose the technology and the shell; decide how
        the domain knowledge will be \emph{represented}; settle the integration with
        existing databases.
  \item \textbf{3. Develop the prototype.} The \textbf{knowledge engineer} acquires the
        domain knowledge from the expert and encodes it as IF--THEN rules. This is the
        long stage.
  \item \textbf{4. Test and refine.} The knowledge engineer runs sample cases; end users
        try the prototype; rules are corrected and added.
        \mk{Iterate 3 and 4 until it performs like the expert.}
  \item \textbf{5. Develop and complete.} Integrate with the real environment, document
        it, train the users.
  \item \textbf{6. Maintain.} Keep the knowledge base current as the domain changes, and
        keep the interfaces working as the surrounding systems evolve.
\end{itemize}
\lead{Answering the scenario version}
\begin{itemize}
  \item \mk{Name the domain, the expert and three real rules.} For
        \yr{\textbf{80 Ch}}'s disease diagnosis: expert $=$ a physician; rules like
        \emph{IF fever $>$ 102\,$^\circ$F AND rash AND joint pain THEN suspect dengue,
        certainty 0.8}; working memory $=$ this patient's symptoms; backward chaining,
        because you are testing a hypothesis.
  \item Then draw the block diagram with the boxes \emph{relabelled for that domain}, and
        walk the six stages in one line each. That is a full 8 marks.
\end{itemize}}{%
\lead{Knowledge acquisition \gap{\color{sub}\footnotesize\tP{2} \yr{\textbf{71 Bh} \m{2+6}, \textbf{74 Bh} \m{2+5}}}}
\begin{itemize}
  \item \textbf{Knowledge acquisition is the process of extracting knowledge from human
        experts and other sources and encoding it into the knowledge base.}
  \item \mk{It is the bottleneck of the whole field.} Experts know
        more than they can say, disagree with each other, and are expensive; and what
        they do say has to be turned into rules by someone who is not an expert.
  \item \textbf{The methods}: structured and unstructured \textbf{interviews};
        \textbf{observation} of the expert working; \textbf{questionnaires};
        \textbf{protocol analysis} (the expert thinks aloud); \textbf{case studies}; and
        \textbf{automatic acquisition}, i.e. machine learning from records (Ch\,6), which
        is what removes the bottleneck when data exists.
  \item \textbf{Sources}: the expert (primary), then secondary and tertiary experts,
        then literature $-$ reports, guidelines, books, manuals $-$ then end users.
  \item \textbf{The four kinds of knowledge to name}: \textbf{declarative} (facts),
        \textbf{procedural} (how to do it), \textbf{episodic} (remembered cases),
        \textbf{meta-knowledge} (knowledge about the knowledge).
  \item \textbf{Example to give}: interviewing a cardiologist about how she reads an ECG,
        recording each judgement as a rule, then replaying old cases to see where the
        rules disagree with her.
\end{itemize}}

%-----------------------------------------------------------------------
\T{7.8 Expert systems: characteristics, advantages, limits, and Vs human experts}

\Q{What are the characteristics of an expert system? Compare it with a human expert. Give its advantages and disadvantages.}{\m{3}\m{4}\m{7}\m{8}\m{2+2+3}}
{\color{sub}\footnotesize \tS{11} \yr{\textbf{72 Ash} \m{8}, \textbf{81 Ch} \m{4+3},
80 Ash \m{8}, 81 Ash \m{4}, \textbf{\texttt{79 Bh}} \m{4}, \textbf{78 Ch} \m{1+4+3},
\textbf{\textit{76 Bh}} \m{5+2+2}, \textbf{77 Ch} \m{6+2}, 82 Ka \m{2+2+3},
\texttt{80 Ba} \m{2+5}, \textbf{\texttt{82 Bh}} \m{2+5}} \gap$\cdot$\gap
\yr{80 Ash} \m{8} asks for the \textbf{comparison alone, with practical examples for each
point}, so the table below \emph{is} that answer}

\sbsr{0.45}{%
\figT{c7_es_vs_human.png}
{\color{sub}\footnotesize Which side wins each point. The top block is where the machine
is better, the bottom block where the person is.}
\lead{Characteristics \gap{\color{sub}\footnotesize(\yr{\texttt{80 Ba}} \m{2}, \yr{82 Ka} \m{2})}}
\begin{itemize}
  \item \textbf{High performance} in its narrow domain, at expert level.
  \item \textbf{Understandable} $-$ takes input and gives output in the user's language.
  \item \textbf{Reliable} and \textbf{highly responsive}.
  \item \textbf{Explains itself}, and separates \textbf{knowledge from the program that
        uses it}, so the knowledge can be edited without touching the code.
  \item Uses \textbf{symbolic} representation and \textbf{heuristics}, not an algorithm.
  \item Handles \textbf{uncertainty}, with certainty factors or probabilities (Ch\,4).
\end{itemize}}{%
\lead{Advantages}
\begin{itemize}
  \item \mk{Permanent}: it does not retire, forget, or die. A human
        expert's knowledge leaves with them.
  \item \textbf{Reproducible} $-$ copy it to a hundred sites at no extra cost, so scarce
        expertise reaches places that have none. This is the argument for
        \yr{\texttt{82 Ba}}'s disaster scenario.
  \item \textbf{Consistent}: same input, same answer, every time. A tired human is not.
  \item \textbf{Fast}, available \textbf{24 hours}, and \textbf{cheaper} than the expert.
  \item Works in \textbf{hazardous environments}; keeps a full audit trail; and is a
        \textbf{training tool}, since it explains its reasoning.
\end{itemize}
\lead{Disadvantages and limitations}
\begin{itemize}
  \item \mk{No common sense.} It has knowledge, not understanding, and
        will give a confidently absurd answer just outside its domain.
  \item \textbf{Narrow}: cannot transfer anything it knows to a neighbouring problem.
  \item \mk{Cannot learn from experience} on its own $-$ every new
        rule has to be put in by a knowledge engineer.
  \item No \textbf{creativity} or intuition, and no ability to recognise that the case in
        front of it is unusual.
  \item \textbf{Expensive and slow to build}, because of the knowledge-acquisition
        bottleneck; and \mk{only as good as the expert it copied},
        errors included.
  \item Cannot take moral responsibility for a decision, which is why medical expert
        systems advise rather than decide.
\end{itemize}}

\sbs{%
\lead{Declarative Vs procedural knowledge \gap{\color{sub}\footnotesize\tF{4} \yr{70 Ma \m{2+6}, 72 Ma \m{3+6}, 73 Ma \m{3+5}, \textit{72 Ma} \m{6}}}}
{\color{sub}\footnotesize all four ask it \emph{with} the expert system architecture, so
the two are one answer.}
\begin{tabularx}{\linewidth}{@{}L{2.0cm}XX@{}}
\hrow \thead{} & \thead{Declarative} & \thead{Procedural} \\
Is & \textbf{knowing what} & \textbf{knowing how} \\
Content & facts about objects, events, situations & the steps to carry a task out \\
Form & statements: logic, semantic nets, frames (Ch\,4, Ch\,5) & rules, programs,
  procedures \\
Stored & separately from how it is used & \textbf{inseparably} from the code using it \\
Changed & easily, one fact at a time & only by rewriting the procedure \\
Reused & by any inference method & only for its own purpose \\
Efficiency & lower: control must be worked out at run time & higher: control is built in \\
Example & \emph{Kathmandu is the capital of Nepal}; \code{dog(Tommy)} & \emph{how to ride
  a bicycle}; a sorting routine \\
\end{tabularx}}{%
\lead{In an expert system, both are present \gap{\color{sub}\footnotesize(this is the half the papers actually mark)}}
\begin{itemize}
  \item The \textbf{knowledge base} is \mk{declarative}: it states what
        is true in the domain and says nothing about how to use it.
  \item The \textbf{inference engine} is \mk{procedural}: it knows how
        to search and fire rules, and nothing about the domain.
  \item \mk{That separation is the whole design}, and the reason a
        shell exists at all: swap the declarative half and the same procedural half
        serves a new domain.
  \item An IF--THEN rule sits on the line, which is worth a sentence: it \emph{states} a
        relationship, so it reads declaratively, but firing it \emph{does} something, so
        it behaves procedurally. That is why production rules are the standard expert
        system representation.
\end{itemize}}

%-----------------------------------------------------------------------
\T{7.9 Natural language processing: NLU, NLG and the five steps}

\Q{What is NLP? Explain the different steps involved, with a block diagram and examples.}{\m{5}\m{6}\m{7}\m{8}\m{9}\m{2+6}\m{3+5}}
{\color{sub}\footnotesize \tS{23} \yr{\textbf{69 Bh} \m{6}, 69 Po \m{8}, 73 Ma \m{2+4},
\textbf{79 Ch} \m{3+5}, \textbf{72 Ash} \m{2+6}, 75 Ba \m{8}, 72 Ma \m{8},
\textbf{\textit{73 Bh}} \m{7}, \textbf{73 Bh} \m{2+7}, \textit{74 Ma} \m{9},
\textbf{\textit{74 Bh}} \m{2+6}, \texttt{82 Ba} \m{2+6}, 79 Jth \m{1+6}, 82 Ka \m{2+5},
\textbf{81 Ch} \m{2+5}, \textbf{\textit{77 Ch}} \m{2+6}, 78 Po \m{3+5}, 79 Ash \m{5+2},
\textbf{75 Bh} \m{6+2+2}, \textbf{\textit{76 Bh}} \m{5+4}, 78 Ba \m{2+6},
\textbf{78 Ch} \m{2+4+2}, \textbf{80 Ch} \m{2+6}}
\gap$\cdot$\gap \mk{the most-asked question in the whole subject, in 23
of the 41 papers}, and it has been set at least once in every year from 2069 to 2082
\gap$\cdot$\gap eight of the twenty-three also want \textbf{NLU and NLG defined}
\gap$\cdot$\gap \yr{\textbf{80 Ch}} asks it as \emph{build a Nepali chatbot}}

\sbsr{0.52}{%
\lead{What NLP is}
\begin{itemize}
  \item \textbf{Natural language processing is the branch of AI that analyses written or
        spoken human language and converts it into a representation a computer can use,
        and back again.}
  \item \textbf{Two halves}, and the papers ask for both by name:
  \begin{itemize}
    \item \textbf{NLU, natural language understanding} $-$ text in, meaning out. Maps the
          input to a useful representation: entities, concepts, keywords, relations,
          sentiment, semantic roles. \mk{This is the hard half.}
    \item \textbf{NLG, natural language generation} $-$ meaning in, text out. Its own
          three stages: \textbf{text planning} (what to say), \textbf{sentence planning}
          (how to phrase and combine it), \textbf{text realisation} (produce grammatical
          output).
  \end{itemize}
  \item \textbf{Input} is text or speech, and \mk{the quality of the
        input bounds everything downstream}: a bad transcription cannot be
        recovered by a good parser.
  \item \textbf{Applications} to name: machine translation, chatbots and virtual
        assistants, sentiment analysis, information retrieval and search, spam
        filtering, text summarisation, question answering, speech recognition, spelling
        and grammar correction, optical character recognition.
\end{itemize}}{%
\figT{c7_nlp_steps.png}
{\color{sub}\footnotesize \mk{Draw these five boxes in this order.} It
is the single most valuable minute of drawing in the subject. (The published figure
prints \emph{Disclosure Integration}; the term is \textbf{Discourse Integration} $-$ a
misprint in the book, not a different step.)}}

\sbsr{0.55}{%
\lead{The five steps, each with the example to give}
\begin{itemize}
  \item \textbf{1. Lexical / morphological analysis.} Break the text into paragraphs,
        sentences and \textbf{tokens}, and analyse the structure of each word.
  \begin{itemize}
    \item \emph{``I love NLP, it's fascinating!''} $\rightarrow$
          \code{[I] [love] [NLP] [,] [it's] [fascinating] [!]}
    \item Includes \textbf{stemming} (rule-based chopping: \emph{closed, closing, closer}
          $\rightarrow$ \emph{clos}) and \textbf{lemmatization} (dictionary-based:
          \emph{better} $\rightarrow$ \emph{good}). Name the difference; it is a
          reliable extra mark.
    \item Morphology: \emph{unhappily} $=$ \emph{un-} $+$ \emph{happy} $+$ \emph{-ly}.
  \end{itemize}
  \item \textbf{2. Syntactic analysis (parsing).} Check the words against the grammar and
        build a \textbf{parse tree} showing how they combine.
  \begin{itemize}
    \item \emph{``The school goes to boy''} is \mk{rejected here} $-$
          every word is a real word, but the arrangement is not English.
    \item Includes \textbf{part-of-speech tagging}: \emph{book} is a noun in
          \emph{the book on the table} and a verb in \emph{to book a flight}.
    \item A worked parse tree is in the companion.
  \end{itemize}
\end{itemize}}{%
\lead{Steps 3 to 5}
\begin{itemize}
  \item \textbf{3. Semantic analysis.} Turn the parse into \textbf{meaning}.
  \begin{itemize}
    \item \textbf{Word sense disambiguation}: \emph{bank} the riverbank or
          \emph{bank} the institution; \emph{Mary had a bat in her office}.
    \item Sentence-level meaning: \emph{``Colourless green ideas sleep furiously''} is
          perfect syntax and \mk{rejected here}, which is the clean
          example of what step 3 does that step 2 cannot.
  \end{itemize}
  \item \textbf{4. Discourse integration.} The meaning of a sentence depends on the
        sentences around it.
  \begin{itemize}
    \item \emph{``The thieves stole the paintings. They were subsequently sold /
          caught / found.''} \mk{One word changes who ``they'' is.}
    \item \emph{``John wanted it''} means nothing without the previous sentence.
  \end{itemize}
  \item \textbf{5. Pragmatic analysis.} What was actually \emph{meant}, using real-world
        knowledge.
  \begin{itemize}
    \item \emph{``Can you pass the salt?''} is a \textbf{request}, not a question about
          your arm. Answering \emph{yes} is grammatically perfect and pragmatically
          wrong.
    \item Covers sarcasm, idiom and implied meaning.
  \end{itemize}
\end{itemize}}

\sbsr{0.55}{%
\lead{The four levels, and the ambiguity at each \gap{\color{sub}\footnotesize\tS{8} \yr{71 Ma \m{2+2+2+2}, \textbf{\textit{71 Bh}} \m{8}, \textit{72 Ma} \m{6}, 78 Ba \m{2+6}, \textbf{\textit{72 Ash}} \m{1+7}, \textbf{\texttt{82 Bh}} \m{3+4}, \textbf{\texttt{80 Bh}} \m{3+5}, \texttt{80 Ba} \m{2+5}}}}
{\color{sub}\footnotesize eight papers ask for the \emph{levels} rather than the steps,
usually with \emph{``give an example of ambiguity at each''}. Same material, different
cut: give one example per level and you have the marks.}
\begin{tabularx}{\linewidth}{@{}L{1.9cm}XX@{}}
\hrow \thead{Level} & \thead{Concerns} & \thead{An ambiguity} \\
\textbf{Phonetic} & sounds $\rightarrow$ words & \emph{``ice cream''} against
  \emph{``I scream''}; \emph{``grey day''} against \emph{``grade A''} \\
\textbf{Lexical} & which sense of a word & \emph{bank}: riverbank or institution \\
\textbf{Syntactic} & which grammatical structure & \emph{``I saw the man with a
  telescope''}: who has the telescope? \\
\textbf{Semantic} & what the sentence means & \emph{``The chicken is ready to eat''}:
  is it the eater or the meal? \\
\textbf{Pragmatic} & what the speaker intended & \emph{``Can you open the window?''}:
  request, not a question \\
\end{tabularx}}{%
\begin{itemize}
  \item \mk{Phonetic is the level every deck skips}, and six papers
        list it. It is the level at which \emph{sound} is turned into \emph{words}, and
        its ambiguity is that one stream of sound splits into words in more than one way.
  \item \textbf{Syntactic and semantic analysis compared} (\yr{\textbf{\texttt{80 Bh}}},
        \yr{\texttt{80 Ba}}, \yr{\textbf{\texttt{82 Bh}}} ask exactly this): syntax asks
        \emph{is this arrangement legal, and what is its structure} and answers with a
        parse tree; semantics asks \emph{what does that structure mean} and answers with
        a meaning representation. A sentence can pass one and fail the other, in both
        directions.
\end{itemize}}

%-----------------------------------------------------------------------
\T{7.10 Issues and challenges in NLP, and machine translation}

\Q{What are the problems associated with NLP? Why is it hard?}{\m{2}\m{3}\m{5}\m{8}\m{1+7}}
{\color{sub}\footnotesize \tS{13} \yr{\textbf{70 Bh} \m{1+7}, 70 Ma \m{2+6},
\textbf{78 Ch} \m{2+4+2}, 82 Ka \m{2+5}, \textbf{81 Ch} \m{2+5},
\textbf{\textit{77 Ch}} \m{2+6}, 78 Po \m{3+5}, 79 Ash \m{5+2}, \textbf{75 Bh} \m{6+2+2},
\textbf{\texttt{79 Bh}} \m{2+5}, 81 Ash \m{8}, \textbf{\texttt{80 Bh}} \m{3+5},
\texttt{80 Ba} \m{2+5}} \gap$\cdot$\gap
\mk{almost always the \m{2} or \m{3} opener to the steps question}
\gap$\cdot$\gap \yr{81 Ash} \m{8} and \yr{\textbf{80 Ch}} \m{2+6} set it as a
\textbf{Nepali-language} problem and want the issues \emph{related to Nepali}}

\sbsr{0.50}{%
\lead{Why natural language is hard}
\begin{itemize}
  \item \textbf{Ambiguity at every level} $-$ \S7.9's table. This is the root cause and
        should be your first sentence.
  \item \textbf{Context dependence}: \emph{``Where's the water?''} means something
        different in a chemistry lab, when you are thirsty, and under a leaking roof.
  \item \textbf{Many ways to say one thing}: \emph{``Ram's birthday is October 11''} and
        \emph{``Ram was born on October 11''}.
  \item \textbf{The language will not stay still}: new words, new senses, slang. No
        lexicon is ever complete.
  \item \textbf{Hidden and figurative meaning}: idiom, metaphor, sarcasm.
  \item \textbf{Pronouns and reference}: \emph{``Ravi went to the supermarket. He found
        his favourite coffee. He paid for it and left.''} What is \emph{it}?
  \item \textbf{Ungrammatical and incomplete input}: \emph{``He rice eats.''} People
        still understand it; a parser does not.
  \item \textbf{Ellipsis and conjunction}: \emph{``Ram and Hari went to a restaurant.
        Ram had coffee, Hari tea.''} $-$ the second verb is simply missing.
  \item \textbf{Speech adds its own}: accents, homophones, and
        \mk{no spaces between spoken words}.
\end{itemize}}{%
\lead{The Nepali version of the answer \gap{\color{sub}\footnotesize(\yr{81 Ash} \m{8}, \yr{\textbf{80 Ch}} \m{6})}}
\begin{itemize}
  \item \mk{Low-resource language}: little annotated text, no large
        treebank, few parallel corpora. Everything below follows from this.
  \item \textbf{Devanagari script}: transliteration and code-mixing (Nepali typed in
        Roman letters), inconsistent Unicode encoding, no standard spelling for many
        loanwords.
  \item \textbf{Morphologically rich}: heavy inflection and agglutination, so one root
        yields dozens of surface forms and a lexicon lookup fails.
  \item \textbf{Honorifics and levels of respect} (\emph{timi / tapai / hajur}) change the
        verb, and carry meaning a translation must preserve.
  \item \textbf{Free word order}, and \textbf{subject-object-verb}, so English-trained
        parsers transfer badly.
  \item \textbf{Dialects} and heavy loan vocabulary from English and Hindi.
  \item \mk{Mitigation, which is the half that carries the marks}:
        build and release an annotated corpus; use transfer learning from a multilingual
        model rather than training from scratch; augment data by back-translation; write
        a morphological analyser for the inflections; standardise the encoding; keep a
        human in the loop for evaluation.
\end{itemize}}

\sbs{%
\lead{Machine translation \gap{\color{sub}\footnotesize(\yr{\textbf{70 Bh}} \m{1+7})}}
\begin{itemize}
  \item \textbf{Machine translation is the automatic conversion of text or speech from
        one natural language into another, preserving meaning.}
  \item \textbf{Three approaches}, in historical order:
  \begin{itemize}
    \item \textbf{Rule based} $-$ dictionaries plus hand-written grammar rules. Accurate
          where the rules cover, brittle everywhere else.
    \item \textbf{Statistical} $-$ learn word and phrase probabilities from a
          \textbf{parallel corpus}. Needs vast aligned data.
    \item \textbf{Neural} $-$ an encoder-decoder network trained end to end. The current
          standard, and the reason translation improved suddenly after 2016.
  \end{itemize}
  \item \textbf{Its challenges}: word-sense ambiguity; idioms that must not be translated
        literally; different \textbf{word order} (English SVO, Nepali SOV); grammatical
        gender and honorifics with no counterpart in the other language; named entities
        and untranslatable culture-specific words; and no parallel corpus for
        low-resource pairs.
\end{itemize}}{%
\lead{Tools NLP needs \gap{\color{sub}\footnotesize(the ``knowledge bases required'', \yr{\textbf{80 Ch}} \m{2})}}
\begin{itemize}
  \item \textbf{Lexicon / dictionary} $-$ the vocabulary, with meanings, pronunciations
        and syntactic categories.
  \item \textbf{Morphological analyser} $-$ splits words into prefix, root and suffix and
        derives their grammatical features.
  \item \textbf{Grammar} $-$ the rules that say which arrangements are legal, used by the
        parser.
  \item \textbf{Semantic and world knowledge} $-$ an ontology or knowledge base (Ch\,5)
        that says what the words refer to and what is plausible.
  \item \textbf{Annotated corpus} $-$ tagged and parsed text to train and to test on.
  \item \textbf{Discourse and pragmatic model} $-$ to resolve pronouns and intent.
\end{itemize}}

%-----------------------------------------------------------------------
\T{7.11 Machine vision}

\Q{What is machine vision? Why is it important in AI? List its steps and applications.}{\m{3}\m{4}\m{2+4}\m{2+5}\m{2+6}\m{3+5}\m{1+5+2}}
{\color{sub}\footnotesize \tS{12} \yr{71 Ma \m{3+5}, \textbf{\textit{72 Ash}} \m{3+5},
80 Ash \m{2+6}, 81 Ash \m{2+4}, \textbf{\texttt{81 Bh}} \m{2+5}, 79 Ash \m{1+5+2},
\textbf{69 Bh} \m{4}, \textbf{80 Ch} \m{4}, \texttt{80 Ba} \m{3}, \texttt{82 Ba} \m{3},
\textbf{\texttt{79 Bh}} \m{2+5}, \textbf{70 Bh} \m{2+6}} \gap$\cdot$\gap
\mk{two papers pair it with back-propagation} (\yr{71 Ma},
\yr{\textbf{\textit{72 Ash}}}) and one with GA (\yr{\textbf{70 Bh}}), so it is often the
\emph{other} half of a neural-network question \gap$\cdot$\gap \yr{80 Ash} wants a
\textbf{scenario} and the steps to solve it}

\sbsr{0.44}{%
\begin{itemize}
  \item \textbf{Machine vision is the field of AI that builds a model of the real world
        from images} $-$ recovering useful information about a three-dimensional scene
        from its two-dimensional projections.
  \item \mk{The hard part in one line}: the world is 3-D and the image
        is 2-D, so information has been thrown away and must be put back using knowledge
        about the objects and the projection geometry.
  \item \textbf{An image} is a 2-D array of intensity values. Each element is a
        \textbf{pixel}. Grey scale $f = f(x,y)$; colour is three such functions,
        $\{R(x,y), G(x,y), B(x,y)\}$.
\end{itemize}}{%
\figT{c7_mv_stages.png}}

\sbsr{0.44}{%
\lead{Why it matters to AI \gap{\color{sub}\footnotesize(\yr{71 Ma} \m{3}, \yr{\textbf{\texttt{79 Bh}}} \m{5})}}
\begin{itemize}
  \item \mk{Vision is how an agent perceives}, and Chapter 1's agent
        model begins with sensors. Without it an agent is limited to whatever a human
        types in.
  \item It is the only practical channel for \textbf{unstructured} information about the
        physical world: no one can hand-code a description of every scene.
  \item It is what makes autonomy possible $-$ robots, vehicles, inspection $-$ and it
        is where neural networks work best, which is why \textbf{back-propagation}
        (\S7.3) and machine vision are asked together.
\end{itemize}}{%
\lead{The stages, in order}
\begin{itemize}
  \item \textbf{1. Image acquisition} $-$ capture the scene with a camera or sensor and
        digitise it. Lighting is a design decision, not an accident.
  \item \textbf{2. Image enhancement} $-$ make features easier to see: contrast
        stretching, sharpening, noise filtering, edge detection.
  \item \textbf{3. Image restoration} $-$ undo a \emph{known} degradation (blur, sensor
        noise) using a model of how it happened. Enhancement is subjective, restoration
        is objective; say so.
  \item \textbf{4. Morphological processing} $-$ operate on shape: erosion, dilation,
        opening, closing, to clean up the regions.
  \item \textbf{5. Segmentation} $-$ \mk{partition the image into
        objects and background}, grouping pixels that share intensity, colour,
        texture or motion. The hardest step, and the one everything after depends on.
  \item \textbf{6. Representation and description} $-$ turn each region into
        \textbf{features}: boundary or area, size, shape, moments, texture.
  \item \textbf{7. Object recognition} $-$ assign a label to each region by
        \textbf{pattern recognition}: statistical (feature vectors into a classifier) or
        structural (decompose into primitives).
  \item \textbf{8. Interpretation} $-$ say what the scene \emph{means}: how many, where,
        moving how, and what to do about it.
  \item \textbf{Colour processing} and \textbf{image compression} sit beside the pipeline
        rather than in it; the block diagram draws them as separate boxes and so should
        you.
\end{itemize}}

\sbs{%
\lead{Applications}
\begin{itemize}
  \item \textbf{Industry}: automated inspection on a production line, reading part
        numbers, robot guidance and pick-and-place, barcode and label verification.
  \item \textbf{Medicine}: X-ray, CT and MRI analysis, tumour detection, cell counting.
  \item \textbf{Transport}: self-driving cars, lane and sign detection, number-plate
        recognition, traffic monitoring.
  \item \textbf{Security}: face recognition, fingerprint and iris biometrics,
        surveillance, intrusion detection.
  \item \textbf{Agriculture}: crop and disease monitoring, fruit grading, yield
        estimation.
  \item \textbf{Documents}: OCR, handwriting recognition, cheque processing.
  \item \textbf{Remote sensing} and \textbf{astronomy}; \textbf{AR and VR}.
\end{itemize}}{%
\lead{Answering the scenario version \gap{\color{sub}\footnotesize(\yr{80 Ash} \m{6}, \yr{81 Ash} \m{4})}}
\begin{itemize}
  \item \mk{Pick a concrete scenario and walk the eight stages
        through it.} That is the answer; a generic list of stages is half marks.
  \item Worked example $-$ \textbf{sorting apples on a conveyor}:
  \begin{itemize}
    \item \textbf{Acquire}: overhead camera, diffuse lighting to kill shadows.
    \item \textbf{Enhance}: median filter to remove sensor speckle.
    \item \textbf{Restore}: deblur for the belt's motion, which is known and constant.
    \item \textbf{Morphology}: close small gaps so each apple is one solid region.
    \item \textbf{Segment}: threshold on colour to separate fruit from belt.
    \item \textbf{Describe}: for each region take area, roundness, mean hue, and the
          fraction of dark pixels.
    \item \textbf{Recognise}: a classifier sorts each into grade A, grade B or reject.
    \item \textbf{Interpret and act}: fire the diverter arm for the right bin.
  \end{itemize}
  \item Say what would break it $-$ overlapping apples defeat segmentation, changing
        daylight defeats the colour threshold $-$ and how you would fix it. That
        sentence is what separates a 6 from a 4.
\end{itemize}}

\hr

%-----------------------------------------------------------------------
\T{7.12 Every Chapter 7 question, grouped by the answer it wants}

{\color{sub}\footnotesize \textbf{Bold} $=$ regular exam, plain $=$ back.
\texttt{Monospace} $=$ CT 710 (BEI), \textit{italic} $=$ CT 78506 (BEX), upright $=$
CT 653 (BCT).}

\begin{itemize}
  \item \textbf{Steps of NLP} \textbf{23 papers} $\rightarrow$ \S7.9. Wordings, all one
        answer: \emph{explain the steps} (\yr{69 Po}, \yr{73 Ma}, \yr{\textbf{79 Ch}},
        \yr{75 Ba}, \yr{72 Ma}, \yr{\textbf{\textit{73 Bh}}});
        \emph{define NLU and NLG, then the steps} (\yr{\textbf{73 Bh}},
        \yr{\textit{74 Ma}}, \yr{\textbf{\textit{74 Bh}}}, \yr{\texttt{82 Ba}},
        \yr{79 Jth}); \emph{the problems, then the steps} (\yr{82 Ka},
        \yr{\textbf{81 Ch}}, \yr{\textbf{\textit{77 Ch}}}, \yr{78 Po});
        \emph{why is NLP hard} (\yr{79 Ash}); \emph{steps and a parse tree for ``A girl
        has a name''} (\yr{\textbf{\textit{76 Bh}}}, worked in the companion);
        \emph{a Nepali chatbot} (\yr{\textbf{80 Ch}}).
  \item \textbf{NLP levels and their ambiguities} \textbf{8 papers} $\rightarrow$
        \S7.9's second half (\yr{71 Ma}, \yr{\textbf{\textit{71 Bh}}},
        \yr{\textit{72 Ma}}, \yr{78 Ba}, \yr{\textbf{\textit{72 Ash}}},
        \yr{\textbf{\texttt{82 Bh}}}, \yr{\textbf{\texttt{80 Bh}}}, \yr{\texttt{80 Ba}}).
        Phonetic is the level to make sure you can define.
  \item \textbf{NLP issues, challenges, machine translation} \textbf{13 papers}
        $\rightarrow$ \S7.10, including both Nepali-language questions
        (\yr{81 Ash}, \yr{\textbf{80 Ch}}).
  \item \textbf{Perceptron and the logic gates} \textbf{17 papers} $\rightarrow$ \S7.2,
        all worked in the companion. \emph{OR} (\yr{72 Ma}); \emph{AND}
        (\yr{\textbf{\textit{73 Bh}}}, \yr{78 Po}, \yr{\textbf{\textit{77 Ch}}});
        \emph{XOR} (\yr{73 Ma}, \yr{\textit{74 Ma}}, \yr{\textbf{\textit{74 Bh}}},
        \yr{\textbf{\texttt{82 Bh}}}, \yr{79 Ash}, \yr{\textbf{77 Ch}});
        \emph{any logic gate} (\yr{79 Jth}); \emph{linear separability}
        (\yr{\textbf{\textit{76 Bh}}}, \yr{\textbf{\texttt{81 Bh}}}); \emph{MLP}
        (\yr{\textit{72 Ma}}, \yr{\textbf{\textit{77 Ch}}}); short note (\yr{69 Po}).
  \item \textbf{Expert system architecture} \textbf{17 papers} $\rightarrow$ \S7.6.
        \textbf{Development stages} \textbf{8 papers} $\rightarrow$ \S7.7.
        \textbf{Advantages, limits, Vs human experts} \textbf{11 papers} $\rightarrow$
        \S7.8. \textbf{Declarative Vs procedural} \textbf{4 papers} $\rightarrow$ \S7.8.
        \textbf{Knowledge acquisition} \textbf{2 papers} $\rightarrow$ \S7.7.
  \item \textbf{Machine vision} \textbf{12 papers} $\rightarrow$ \S7.11.
  \item \textbf{Kohonen / SOM} \textbf{9 papers} $\rightarrow$ \S7.5.
        \textbf{Hopfield} \textbf{8 papers} $\rightarrow$ \S7.4, with the weight matrix
        and recall worked in the companion for \yr{\textbf{\texttt{82 Bh}}}.
  \item \textbf{Back-propagation} \textbf{8 papers} $\rightarrow$ \S7.3, one full step
        worked in the companion. \textbf{Hebbian learning} \textbf{4 papers}
        $\rightarrow$ \S7.3, the AND net worked in the companion.
  \item \textbf{The neuron model, McCulloch-Pitts, activation function}
        \textbf{8 papers} $\rightarrow$ \S7.1.
  \item \textbf{Forward propagation through the printed network} \tP{2}
        \yr{82 Ka \m{8}, \textbf{81 Ch} \m{8}} $-$ the only genuine numerical in the
        chapter, and \mk{the same question printed twice with the same
        figure}. Worked in the companion, including why the ReLU layer's weights
        cannot be updated at all for the given input.
  \item \textbf{Short notes that live here}: \emph{learning in a neural network}
        (\yr{69 Po}, \yr{\textbf{74 Bh}}, \yr{78 Ba}, \yr{\textbf{\textit{77 Ch}}},
        \yr{81 Ash}), \emph{learning rate} and \emph{learning based on statistics}
        (\yr{\textit{72 Ma}}), \emph{human brain Vs neural network}
        (\yr{\textbf{69 Bh}}).
  \item \textbf{Adaline} $-$ on the syllabus, \mk{never examined in 41
        papers}. Four lines in \S7.3 and no more.
  \item \textbf{Pattern recognition} $-$ in the textbook's \S7.4 but not on the CT 710
        syllabus and never asked on its own. It appears here only as step 7 of the
        machine-vision pipeline.
\end{itemize}

\lead{Source notes}
\begin{itemize}
  \item Spine: BA Sir \emph{CH-07 Applications of AI} (the fullest single source in
        \emph{Notes\textbackslash} for this chapter), PS Sir \emph{7.1 Neural Network},
        \emph{7.2 Expert Systems} and \emph{7.3 NLP and Machine Vision}, and
        \emph{Insights on Artificial Intelligence} ch\,7 for the worked method.
  \item \mk{One correction carried into the figure caption}: the
        published NLP figure prints \emph{Disclosure Integration}; the step is
        \textbf{Discourse Integration}. The same misprint runs through several decks.
  \item \mk{Phonetic-level ambiguity is \added{}}. Six papers ask for
        the four levels of analysis and no deck in \emph{Notes\textbackslash} explains
        the phonetic one; the examples here were researched and written.
  \item \mk{The Nepali-language answers are \added{}} for
        \yr{81 Ash} and \yr{\textbf{80 Ch}}: no source covers them, and both questions
        are worth 6 or 8 marks.
  \item \textbf{Hebbian learning} is not on the syllabus under any name but is asked in
        four papers, so it is written in full in \S7.3 and worked in the companion.
\end{itemize}
